\begin{frame}
  \titlepage
\end{frame}

\begin{frame}{15-Minute Roadmap and Rubric Match}
\begin{enumerate}
    \item Problem, research question, and contribution.
    \item New clinical evidence: COVID test routing and pulse-ox bias.
    \item Updated domain model and architecture.
    \item Best analysis artifacts from \texttt{presentation/analysis\_artifacts}.
    \item Results, code-review targets, demo path, limitations, and next-semester plan.
\end{enumerate}
\begin{block}{Presentation requirement alignment}
This deck includes background, motivation, updated architecture, code-review targets, current results, graph-backed clinical validation, demo plan, and next steps.
\end{block}
\end{frame}

\begin{frame}{Project Claim}
\begin{block}{Scientific claim}
Context quality is a fairness variable: when sequential decision systems learn from incomplete or unequally reliable context, they can repeatedly route scarce resources away from the groups or flows that need them most.
\end{block}
\begin{columns}[T]
\begin{column}{0.48\textwidth}
\textbf{Clinical routing}
\begin{itemize}
    \item tests, retests, and triage escalation
    \item under-measured or mis-measured groups
    \item harm as missed diagnosis or delayed care
\end{itemize}
\end{column}
\begin{column}{0.48\textwidth}
\textbf{Quantum routing}
\begin{itemize}
    \item route selection and qubit allocation
    \item stale or incomplete link context
    \item harm as repeated routing failure
\end{itemize}
\end{column}
\end{columns}
\end{frame}

\begin{frame}{Why Clinical and Quantum Belong in the Same Study}
\begin{center}
\begin{tabular}{p{0.25\textwidth}p{0.30\textwidth}p{0.31\textwidth}}
\toprule
\textbf{Shared feature} & \textbf{Clinical testbed} & \textbf{Quantum testbed} \\
\midrule
Scarce resource & diagnostic attention, retests, escalation & route capacity, qubits, entanglement effort \\
Partial feedback & chosen test or escalation only & chosen path/allocation only \\
Context failure & missing risk or calibration context & stale topology or link context \\
Fairness metric & group/site/cohort disparity & flow/class/service disparity \\
\bottomrule
\end{tabular}
\end{center}
\vspace{0.5em}
\alert{Research goal:} compare \mab, \cmab, and \icmab policies under one utility/fairness reporting surface.
\end{frame}

\begin{frame}{New Clinical Foundation: Two COVID Cases}
\begin{columns}[T]
\begin{column}{0.48\textwidth}
\textbf{Case 1: Testing allocation}
\begin{itemize}
    \item high-SVI Hispanic and Black communities had high burden
    \item testing access was often lower or delayed
    \item fair CMAB routes tests by burden and vulnerability
\end{itemize}
\end{column}
\begin{column}{0.48\textwidth}
\textbf{Case 2: Pulse oximetry}
\begin{itemize}
    \item calibration context was incomplete across skin tone
    \item same SpO$_2$ threshold produced unequal reliability
    \item EQUITAS mediation targets ambiguous decisions
\end{itemize}
\end{column}
\end{columns}
\vspace{0.4em}
\scriptsize Sources include CDC COVID disparities, Sjoding et al. 2020, Fawzy et al. 2022/2023, FDA pulse oximetry safety communication, and the new COVID clinical foundation report.
\end{frame}
